\documentclass[11pt]{article}

\usepackage[review]{acl}

\usepackage{times}
\usepackage{latexsym}
\usepackage[T1]{fontenc}
\usepackage[utf8]{inputenc}
\usepackage{microtype}
\usepackage{inconsolata}
\usepackage{graphicx}
\usepackage{booktabs}

\title{Reproducing and Extending KELE: A Multi-Agent Framework for Structured Socratic Teaching with LLMs}

\author{Maximilian Khan \\
  Santa Clara University \\
  \texttt{mjkhan@scu.edu} \\\And
  Ulises Chavarria \\
  Santa Clara University \\
  \texttt{uchavarria@scu.edu} \\}

\begin{document}
\maketitle

\begin{abstract}
We reproduce and extend KELE \citep{peng2025kele}, a multi-agent framework for structured Socratic teaching with large language models. KELE decomposes Socratic dialogue into five progressive teaching stages governed by a rule system called SocRule, coordinated through a ``consultant--teacher'' mechanism in which one LLM plans instruction while another executes it. We reimplement the full pipeline, serve SocratTeachLLM locally via vLLM, and evaluate on the 681-dialogue test split of SocratDataset using ROUGE, BLEU, and stage-level state accuracy. Our baseline reproduction achieves ROUGE-2 of 26.04 and BLEU-4 of 19.60, confirming the core findings while revealing a gap in mid-stage state transitions (stages c/d accuracy below 6\%) that motivates our planned improvement: replacing the LLM-based consultant with a lightweight learned classifier trained on the 42,000-turn dataset labels. Code and data are available at \url{https://github.com/ulises-c/csen-346} and \url{https://huggingface.co/yuanpan/SocratTeachLLM}.
\end{abstract}

%% ============================================================
%% SECTION 1 — INTRODUCTION
%% ============================================================
\section{Introduction}

Dialogue plays a central role in shaping learning outcomes. Traditional knowledge-imparting instruction places students in a passive role, delivering answers directly without engaging deeper reasoning \citep{johnston1994teaching}. Socratic teaching offers an alternative: through structured questioning, learners are guided to actively construct knowledge and discover answers themselves \citep{seeskin1987socratic, chang1998socratic}. Research consistently shows that Socratic dialogue promotes cognitive development and deeper understanding \citep{knezic2010socratic}, yet its reliance on highly skilled instructors makes it difficult to scale in practice.

Recent advances in large language models (LLMs) have opened a path toward automating Socratic instruction. LLMs are increasingly applied in educational settings---student Q\&A, lesson planning, and adaptive tutoring \citep{kasneci2023chatgpt, gan2023llmedu}---and several systems have explored Socratic-style dialogue generation, including prompt-based approaches \citep{zhang2024spl, alhossami2024socratic} and fine-tuned models such as SocraticLM \citep{liu2024socraticLM} and EduChat \citep{dan2023educhat}. However, these efforts generally focus on surface-level heuristic questioning while neglecting the structured, multi-stage progression that characterizes effective Socratic pedagogy.

\citet{peng2025kele} address this gap with KELE (Knowledge-Enlightened Learning Enhanced by LLMs), a multi-agent framework that decomposes Socratic teaching into five hierarchical stages---student questioning, concept probing, inductive reasoning, rule construction, and teacher summary---governed by a formal rule system called SocRule containing 34 teaching strategies. The framework employs a ``consultant--teacher'' mechanism: a \emph{consultant} agent evaluates student cognitive state and selects the appropriate teaching strategy, while a \emph{teacher} agent generates the actual dialogue response. To power the teacher agent, the authors fine-tune GLM4-9B \citep{glm2024chatglm} using LoRA \citep{hu2022lora} on SocratDataset, a corpus of 6,803 dialogues with over 42,000 annotated turns, producing SocratTeachLLM---a model that outperforms GPT-4o across all nine Socratic teaching evaluation dimensions despite being significantly smaller.

In this work, we reproduce and extend the KELE framework with three goals:

\begin{enumerate}
    \item \textbf{Reproduction.} We reimplement the consultant--teacher pipeline, serve SocratTeachLLM locally via vLLM \citep{kwon2023vllm}, and use GPT-4o \citep{openai2024gpt4o} as the consultant to evaluate on the same 681-dialogue test split. We report ROUGE \citep{lin2004rouge}, BLEU \citep{papineni2002bleu}, and per-stage state accuracy to verify the original claims.
    \item \textbf{Diagnosis.} We analyze where the reproduction matches or diverges from the published results, with particular attention to the low mid-stage state accuracy (stages c and d below 6\%) that emerges from our evaluation.
    \item \textbf{Improvement.} We propose replacing the LLM-based consultant with a lightweight supervised classifier trained on SocratDataset's 42,000 state-labeled turns, targeting the system's weakest link: the consultant's ability to correctly identify student cognitive states during the complex middle stages of Socratic dialogue.
\end{enumerate}

%% ============================================================
%% SECTION 2 — RELATED WORK
%% ============================================================
\section{Related Work}

\subsection{Large Language Models in Education}

LLMs are rapidly being integrated across educational applications, from student Q\&A and personalized tutoring to lesson planning and intelligent feedback \citep{kasneci2023chatgpt, gan2023llmedu}. InstructGPT-style alignment through reinforcement learning from human feedback has enabled conversational agents that can adapt explanations to student needs \citep{chung2024scaling}. RetLLM-E \citep{mitra2024retllm} retrieves relevant course material to ground LLM responses for student forum questions, while LessonPlanner \citep{fan2024lessonplanner} assists novice teachers in generating pedagogically sound lesson plans. The CLASS framework \citep{sonkar2023class} builds intelligent tutoring systems grounded in learning science principles, guiding students step-by-step through interactive dialogue. These systems demonstrate the growing capacity of LLMs to serve as educational tools, though most operate in a single-agent paradigm without explicit pedagogical structure.

\subsection{Technology-Enhanced Socratic Teaching}

The Socratic method---teaching through guided questioning rather than direct instruction---has a long history in educational research \citep{seeskin1987socratic, chang1998socratic}. Technology-enhanced implementations range from early rule-based dialogue systems \citep{graesser2004autotutor} to modern LLM-powered approaches. \citet{alhossami2024socratic} demonstrate that LLMs can employ Socratic questioning to help students debug code, while \citet{shridhar2022mathword} show automatic generation of Socratic sub-questions for math word problems. The SPL system \citep{zhang2024spl} leverages GPT-4 for Socratic dialogue in a playground-style interface.

On the fine-tuning side, EduChat \citep{dan2023educhat} and SocraticLM \citep{liu2024socraticLM} train models on Socratic-style datasets to internalize heuristic questioning patterns, while \citet{ding2024socraticmath} apply the Socratic method specifically to conversational mathematics teaching. MathDial \citep{macina2023mathdial} contributes a tutoring dialogue dataset grounded in math reasoning with rich pedagogical annotations.

However, as \citet{peng2025kele} observe, these approaches primarily focus on generating heuristic questions at the individual turn level without enforcing coherent multi-stage teaching progression. The absence of structured rules can lead to arbitrary guidance, abrupt topic shifts, and failure to consolidate learning. KELE addresses this through its five-stage SocRule system and dual-agent architecture, making it---to our knowledge---the first framework to systematically implement structured Socratic teaching with LLMs.

\subsection{Multi-Agent Collaborative Systems}

Multi-agent architectures have emerged as an effective paradigm for complex LLM tasks that benefit from separation of concerns \citep{ge2023openagi}. In the educational domain, the key insight of KELE's consultant--teacher design is that teaching \emph{planning} (what stage and strategy to use) and teaching \emph{execution} (generating the actual dialogue) are fundamentally different capabilities that benefit from distinct agents \citep{peng2025kele}. This mirrors real-world educational practice where teaching consultants advise on instructional strategy while classroom teachers handle direct student interaction \citep{nevin2009collaborative}. Single-agent approaches struggle with both responsibilities simultaneously, often producing responses that are individually reasonable but lack coherent pedagogical progression \citep{an2024llmcontext}. Our work investigates whether the consultant role---essentially a state classification task---can be more effectively handled by a specialized classifier rather than a general-purpose LLM.

%% ============================================================
%% SECTION 3 — DATASET & METHODOLOGY (placeholder)
%% ============================================================
\section{Dataset and Methodology}

% To be completed for the Apr 23 submission.

%% ============================================================
%% SECTION 4 — EVALUATION & RESULTS (placeholder)
%% ============================================================
\section{Evaluation and Results}

% To be completed for the May 5 submission.

%% ============================================================
%% SECTION 5 — CONCLUSION (placeholder)
%% ============================================================
\section{Conclusion}

% To be completed for the May 14 submission.

%% ============================================================
%% LIMITATIONS (placeholder)
%% ============================================================
\section*{Limitations}

% To be completed for the May 14 submission.

%% ============================================================
%% ETHICS (placeholder)
%% ============================================================
\section*{Ethics Statement}

% To be completed for the May 14 submission.

\bibliography{custom}

\end{document}
